%% kapa: exam=1 round=26 year=2015 period=1 form=A subject=civil_law subject_ko=민법 seq=26 number=26
%% source: Q-Net 감정평가사 기출문제 (cst00309) artlSeq=5138001 file=extracted/1008/5138001/제26회_감정평가사 1차_기출/제26회_감정평가사 1차_기출/1교시_감평_A형.pdf
%% conversion: kordoc 4.24.51 -> tools/gen1.py (auto); review=proofread-sonnet; flags=none
\begin{question}{26}
소유권 취득에 관한 설명으로 옳지 않은 것은? (다툼이 있으면 판례에 따름)
\begin{choices}
  \item 무주의 부동산도 선점의 대상이 된다.
  \item 부동산 매수인이 매도인의 부동산 처분권한을 조사했더라면 그 처분권한이 없음을 알 수 있었음에도 이를 조사하지 않은 경우, 매수인의 등기부취득시효는 완성되지 않는다.
  \item 부합한 동산의 주종을 구별할 수 없는 때에는 동산의 소유자는 부합 당시의 가액의 비율로 합성물을 공유한다.
  \item 타인의 토지 기타 물건으로부터 발견된, 문화재가 아닌 매장물은 법률이 정한 바에 의하여 공고한 후 1년 내에 그 소유자가 권리를 주장하지 아니하면 그 토지 기타 물건의 소유자와 발견자가 절반하여 취득한다.
  \item 타인의 권원에 의하여 부동산에 부합된 물건이 독립한 권리의 객체성을 상실하고 부동산의 구성부분이 된 경우, 그 부합물의 소유권은 부동산의 소유자에게 귀속된다.
\end{choices}
\end{question}
\answer{1}
